\documentclass[12pt]{article}
\usepackage[utf8]{inputenc}
\usepackage[T1]{fontenc}
\usepackage{hyperref}
\title{Komviri Language Resources for K-12 Curriculum Development}
\author{Compiled by Claude Sonnet 5.5 (claude-sonnet-5-5) for Andr\'as Kornai}
\date{}
\begin{document}
\maketitle

\section{Introduction}
Komviri (Kamviri, Kamkata-viri; ISO 639-3 xvi) is a Nuristani language of the Indo-Iranian branch of Indo-European. It is spoken in Afghanistan and Pakistan. Some sources treat it as a dialect of Katë. Speaker estimates in the search results range from about 5,000--10,000 to about 20,000. No school materials were found, and no published full grammar exists.

\section{Primary Educational Resources}
No primers, textbooks or school programmes were found.

\section{Linguistic and Reference Materials}
\begin{itemize}
\item Richard Strand's Nuristan materials, including the Cologne workshop page, which is the source for his Kamviri documentation: \url{https://nuristan.info/CologneWorkshop/CologneWorkshop.html}
\item Glottolog entry for Ghul\^am\=ull\^ah, Q\^azi (1966), a manuscript grammar of a Nuristani language: \url{https://glottolog.org/resource/reference/id/571727}
\item IECoR (Indo-European Cognate Relationships) database, contributors page: \url{https://iecor.clld.org/contributors/27}
\item Wikipedia, Kamviri dialect: \url{https://EN.wikipedia.com/wiki/Kamviri}
\item Wikipedia, Kat\"e language: \url{https://en.wikipedia.org/wiki/Kat%C3%AB_language}
\item Wikipedia, Nuristani languages: \url{https://en.wikipedia.org/wiki/Nuristani_languages}
\end{itemize}
Strand's ``K\^amv'iri Lexicon'' and ``The Sound System of K\^amv'iri'' (1997) are mentioned in the search summaries as being online. I did not see their direct URLs, so I do not cite them.

\section{Religious and Community Resources}
None found.

\section{Digital and Audio Resources}
No audio resources were found beyond what the Strand site may hold.

\section{Curriculum Development Context}
Little educational material exists, so a curriculum would have to be built from scratch. The first step would be community consultation. Strand's documentation and the manuscript grammar are the available linguistic starting points.
\end{document}